\section{Giới thiệu}

Sự phát triển mạnh mẽ của Internet vạn vật (Internet of Things --- IoT) đặt ra yêu cầu ngày càng cao đối với các hệ thống nhúng về khả năng xử lý đa tác vụ, độ tin cậy theo thời gian thực và khả năng tích hợp với hạ tầng đám mây.
Báo cáo này trình bày quá trình thiết kế, triển khai và đánh giá một hệ thống nhúng thời gian thực hoàn chỉnh trên nền tảng vi điều khiển ESP32-S3 (bo mạch Yolo UNO), sử dụng hệ điều hành thời gian thực FreeRTOS~\cite{freertos} làm nền tảng điều phối tác vụ.
Dự án được xây dựng dựa trên mã nguồn gốc YoloUNO PlatformIO RTOS\_Project~\cite{yolouno} và mở rộng với các chức năng khác biệt ít nhất 30\% so với mã nguồn gốc, bao gồm kiến trúc đồng bộ hóa không có biến toàn cục, giao diện giám sát Web theo thời gian thực, suy luận máy học tại biên (TinyML) và tích hợp nền tảng đám mây CoreIOT~\cite{coreiot}.

\subsection{Mục tiêu của hệ thống}

Mục tiêu tổng quát của đề tài là phát triển một hệ thống giám sát môi trường và điều khiển thiết bị ngoại vi hoàn chỉnh trên vi điều khiển ESP32-S3, đảm bảo toàn bộ việc chia sẻ dữ liệu giữa các tác vụ được thực hiện thông qua các cơ chế đồng bộ hóa của FreeRTOS (Mutex và Binary Semaphore) mà không sử dụng bất kỳ biến toàn cục nào.
Cụ thể, hệ thống cần hoàn thành sáu nhiệm vụ độc lập nhưng được tích hợp chặt chẽ với nhau như sau:

\begin{itemize}
    \item \textbf{Nhiệm vụ 1: Đèn LED đơn nhấp nháy theo điều kiện nhiệt độ.}
    Xác định lại chế độ nhấp nháy của đèn LED onboard (GPIO~38) để phản ứng với ít nhất ba khoảng nhiệt độ khác nhau, mỗi khoảng tương ứng với một tần số nhấp nháy riêng biệt.
    Cơ chế Binary Semaphore (\texttt{xLedSemaphore}) được sử dụng để đồng bộ hóa giữa tác vụ đọc cảm biến và tác vụ điều khiển LED, đảm bảo tính nhất quán của dữ liệu nhiệt độ trong môi trường đa tác vụ.

    \item \textbf{Nhiệm vụ 2: Điều khiển đèn LED NeoPixel theo độ ẩm.}
    Thiết kế lại các mẫu màu của đèn NeoPixel WS2812B onboard (GPIO~18) để biểu thị ít nhất ba mức độ ẩm khác nhau, với sự tương quan rõ ràng giữa từng khoảng giá trị độ ẩm và màu sắc hiển thị tương ứng.
    Kỹ thuật Binary Semaphore (\texttt{xNeoSemaphore}) được áp dụng để đồng bộ hóa quá trình cập nhật màu sắc giữa tác vụ cảm biến và tác vụ điều khiển NeoPixel.

    \item \textbf{Nhiệm vụ 3: Giám sát nhiệt độ và độ ẩm qua màn hình LCD.}
    Định nghĩa các điều kiện ngưỡng để tạo và giải phóng Mutex bảo vệ tài nguyên màn hình LCD 1602 I2C, hiển thị ít nhất ba trạng thái cảnh báo khác biệt (Normal, Warning, Critical) dựa trên giá trị đọc từ cảm biến DHT20~\cite{dht20}.
    Nhiệm vụ này cũng yêu cầu loại bỏ hoàn toàn tất cả biến toàn cục trong toàn bộ dự án, thay thế bằng cấu trúc dữ liệu dùng chung \texttt{SystemData} được bảo vệ bởi Mutex (\texttt{xMutex}).

    \item \textbf{Nhiệm vụ 4: Máy chủ Web ở chế độ điểm truy cập.}
    Thiết kế lại giao diện máy chủ web nhúng (\texttt{ESPAsyncWebServer}~\cite{espassyncwebserver}) để nâng cao khả năng sử dụng, kết hợp giao thức WebSocket (\texttt{/ws}) cho phép cập nhật dữ liệu cảm biến theo thời gian thực mà không cần tải lại trang.
    Giao diện phải bao gồm ít nhất hai nút điều khiển có nhãn rõ ràng để điều khiển độc lập hai thiết bị ngoại vi; toàn bộ tài nguyên giao diện được lưu trữ trên phân vùng LittleFS~\cite{littlefs} của bộ nhớ Flash.

    \item \textbf{Nhiệm vụ 5: Triển khai TinyML và đánh giá độ chính xác.}
    Mô tả tập dữ liệu huấn luyện bao gồm quá trình thu thập và gán nhãn dữ liệu, huấn luyện mạng nơ-ron phân loại và chuyển đổi sang định dạng C-array để nhúng trực tiếp vào chip.
    Mô hình được thực thi bằng thư viện TensorFlow Lite Micro~\cite{tflitemicro} trên Tensor Arena tĩnh 8~kB, với quá trình đo lường và đánh giá độ chính xác nhận dạng (Accuracy, Precision, Recall) trên tập kiểm thử, kèm thảo luận về hiệu suất và giới hạn của mô hình khi chạy trên phần cứng thực.

    \item \textbf{Nhiệm vụ 6: Xuất bản dữ liệu lên máy chủ đám mây CoreIOT.}
    Cấu hình ESP32-S3 ở chế độ Trạm (STA Mode) để thiết lập kết nối MQTT đến máy chủ CoreIOT~\cite{coreiot}, định kỳ xuất bản dữ liệu viễn trắc nhiệt độ và độ ẩm cùng thuộc tính định danh thiết bị (địa chỉ MAC).
    Thiết bị sử dụng mã xác thực (Device Token) hợp lệ khớp với thiết bị đã đăng ký trên nền tảng, tuân theo mẫu giải pháp chính thức của CoreIOT và thư viện ThingsBoard MQTT~\cite{thingsboard_mqtt}.
\end{itemize}

\subsection{Phạm vi và giới hạn}

Dự án được triển khai toàn bộ trên nền tảng vi điều khiển ESP32-S3 (bo mạch Yolo UNO) bằng ngôn ngữ C/C++ thông qua plugin PlatformIO~\cite{platformio}. Trong giới hạn của đề tài, cấu hình phần cứng của bo mạch được cố định ở chế độ 301B9 và 812H6.

Về mặt phần cứng, hệ thống sử dụng cảm biến DHT20~\cite{dht20} kết nối qua giao thức I2C (SDA: GPIO~11, SCL: GPIO~12) để thu thập nhiệt độ và độ ẩm, cùng các thiết bị ngoại vi onboard gồm LED đơn (GPIO~38), NeoPixel WS2812B~\cite{ws2812b} (GPIO~18) và màn hình LCD 1602 I2C.
Về mặt kết nối, hệ thống giới hạn ở giao tiếp Wi-Fi 2.4~GHz với một mạng nội bộ duy nhất.
Quá trình huấn luyện mô hình TinyML sử dụng tập dữ liệu có quy mô nhỏ, phù hợp với dung lượng RAM tĩnh 8~kB Tensor Arena của ESP32-S3, và nền tảng CoreIOT được sử dụng như giải pháp lưu trữ đám mây duy nhất thông qua thư viện ThingsBoard MQTT~\cite{thingsboard_mqtt}.
Việc kiểm thử hệ thống được thực hiện trong môi trường phòng thí nghiệm, tập trung vào tính đúng đắn của logic đồng bộ hóa đa tác vụ RTOS thay vì tối ưu hóa cho điều kiện vận hành khắc nghiệt ngoài trời.
